\documentclass[opre]{informs3}
\OneAndAHalfSpacedXI

\usepackage[utf8]{inputenc}
\usepackage{amsmath,amssymb,mathrsfs}
\usepackage{bm,bbm,booktabs,graphicx}
\usepackage{natbib}
\bibpunct[, ]{(}{)}{,}{a}{}{,}%
\def\bibfont{\small}%
\def\bibsep{\smallskipamount}%
\def\bibhang{24pt}%
\def\newblock{\ }%
\def\BIBand{and}%
\usepackage[hyperfootnotes=false,colorlinks=true,linkcolor=black,urlcolor=black,citecolor=blue]{hyperref}
\usepackage{cleveref}

\TheoremsNumberedThrough
\EquationsNumberedThrough
\allowdisplaybreaks

\newcommand{\E}{\mathbb{E}}
\newcommand{\Pb}{\mathbb{P}}
\newcommand{\Natu}{\mathbb{N}}
\newcommand{\Ind}[1]{\mathbbm{1}\{#1\}}
\newcommand{\State}{\mathcal{S}}
\newcommand{\Xscr}{\mathcal{X}}

\begin{document}

\RUNAUTHOR{Natalia Trigo}
\RUNTITLE{Finite-Horizon Dynamic Reputation Model}

\TITLE{A Finite-Horizon Dynamic Reputation Model with Hidden Product Choice}

\ARTICLEAUTHORS{
\AUTHOR{Natalia Trigo}
\AFF{Anderson UCLA, \EMAIL{natalia.trigo.phd@anderson.ucla.edu}}
}

\ABSTRACT{
We study a finite-horizon reputation problem in which a seller privately chooses
between a low-cost product and a high-quality product while a user learns only
from realized binary outcomes. The user treats the seller as a single Bernoulli
arm and applies Thompson sampling to decide whether to engage with the seller or
with a known outside option. The seller chooses hidden products to maximize total
expected profit over a finite horizon. The resulting dynamic program is indexed
by calendar time and by the user's posterior sufficient statistics. The optimal
product choice is governed by the marginal continuation value of a success
relative to a failure, which gives a simple threshold condition for using the
costly product. We also derive a deterministic fluid approximation that replaces
random posterior counts by their expected drifts. The fluid formulation
preserves the same reputation-investment threshold and provides a transparent
long-run benchmark for the quality mix needed to sustain demand.
}

\KEYWORDS{dynamic reputation; hidden actions; finite-horizon dynamic programming; fluid approximation; Thompson sampling}

\maketitle

\section{Model Formulation}

Consider a discrete-time horizon $t=1,\ldots,T$. The environment consists of a
single user and two sellers, denoted by $A$ and $B$. At each period $t$, the user
selects a seller, denoted by $y_t\in\{A,B\}$. The user's realized utility is a
binary random variable $u_t\in\{0,1\}$.

Seller B offers a transparent outside option. If the user selects Seller B, the
realized utility is Bernoulli with known success probability $p_0\in(0,1)$:
\begin{equation}
    \Pb(u_t=1\mid y_t=B)=p_0.
    \label{eq:seller-b-success}
\end{equation}

Seller A has access to two products, indexed by $\Xscr=\{1,2\}$. When the user
selects Seller A, the seller privately chooses a product $x_t\in\Xscr$. Product
$i\in\Xscr$ generates a successful outcome with probability $p_i$ and imposes
cost $c_i$ on Seller A:
\begin{equation}
    \Pb(u_t=1\mid y_t=A,x_t=i)=p_i.
    \label{eq:seller-a-success}
\end{equation}
We assume
\begin{equation}
    0<p_1<p_2\leq 1,\qquad c_1<c_2,
    \label{eq:product-ordering}
\end{equation}
so product 2 is more likely to satisfy the user but is also more costly. Whenever
Seller A is chosen, it receives a fixed revenue $R>0$. Seller A's period profit
is
\begin{equation}
    \pi_t^A=
    \begin{cases}
        R-c_{x_t} & \text{if } y_t=A,\\
        0 & \text{if } y_t=B.
    \end{cases}
    \label{eq:period-profit}
\end{equation}

\section{User Learning and Demand}

The product chosen by Seller A is hidden from the user. The product affects the
distribution of realized utility, but the user observes only the seller selected
and the binary realized outcome.

\begin{assumption}[User's Misspecified Belief]
The user observes $u_t$ after choosing Seller A but does not observe the product
$x_t\in\Xscr$. The user therefore treats Seller A as a single product with an
unknown Bernoulli success probability $\theta\in[0,1]$.
\end{assumption}

The user learns about $\theta$ through Bayesian updating. Let $S_t$ and $F_t$
denote the cumulative number of successes and failures generated by Seller A up
to the beginning of period $t$:
\begin{align}
    S_t &= \sum_{\tau=1}^{t-1}\Ind{y_\tau=A,u_\tau=1},\\
    F_t &= \sum_{\tau=1}^{t-1}\Ind{y_\tau=A,u_\tau=0}.
\end{align}
The user starts from the uniform prior $\theta\sim\operatorname{Beta}(1,1)$.
Given state $(S_t,F_t)$, the posterior belief is
\begin{equation}
    \theta\mid S_t,F_t\sim\operatorname{Beta}(1+S_t,1+F_t).
    \label{eq:posterior}
\end{equation}
The corresponding posterior mean is
\begin{equation}
    m(S_t,F_t)=\frac{S_t+1}{S_t+F_t+2}.
    \label{eq:posterior-mean}
\end{equation}

The user follows a Thompson sampling heuristic. At the beginning of period $t$,
the user samples
\begin{equation}
    \tilde{\theta}_t\sim\operatorname{Beta}(1+S_t,1+F_t).
\end{equation}
Since the outside option has known success probability $p_0$, the user chooses
Seller A if and only if the sampled index exceeds $p_0$:
\begin{equation}
    y_t=
    \begin{cases}
        A & \text{if } \tilde{\theta}_t\geq p_0,\\
        B & \text{if } \tilde{\theta}_t<p_0.
    \end{cases}
    \label{eq:user-policy}
\end{equation}

Seller A observes the user's choices and realized outcomes, and can therefore
track $(S_t,F_t)$. For a state $(S,F)$, define the demand probability
\begin{equation}
    D(S,F)
    \triangleq
    \Pb(\tilde{\theta}_t\geq p_0\mid S_t=S,F_t=F)
    =
    \int_{p_0}^{1}
    \frac{\theta^S(1-\theta)^F}{B(1+S,1+F)}\,d\theta,
    \label{eq:demand-probability}
\end{equation}
where $B(\cdot,\cdot)$ is the Beta function. This probability is the chance that
Seller A has an opportunity to earn revenue and update the user's belief in the
current period.

\section{Finite-Horizon Seller Problem}

Seller A chooses hidden products to maximize total expected profit over the
finite horizon. Let $V_t(S,F)$ denote the optimal expected continuation profit
from periods $t,\ldots,T$ when the user's belief state at the beginning of period
$t$ is $(S,F)$. The terminal condition is
\begin{equation}
    V_{T+1}(S,F)=0
    \qquad\text{for all reachable }(S,F).
    \label{eq:terminal-condition}
\end{equation}

At state $(S,F)$ in period $t$, the user chooses Seller A with probability
$D(S,F)$. If the user chooses Seller B, Seller A earns zero profit and the belief
state remains $(S,F)$. If the user chooses Seller A and the seller uses product
$x\in\{1,2\}$, Seller A earns $R-c_x$ and the next state is $(S+1,F)$ after a
success and $(S,F+1)$ after a failure. The finite-horizon Bellman equation is
\begin{equation}
\begin{aligned}
    V_t(S,F)
    =
    &[1-D(S,F)]V_{t+1}(S,F)\\
    &+
    D(S,F)
    \max_{x\in\{1,2\}}
    \left\{
        R-c_x
        +p_xV_{t+1}(S+1,F)
        +(1-p_x)V_{t+1}(S,F+1)
    \right\}.
\end{aligned}
\label{eq:finite-bellman}
\end{equation}

For each product $x$, define the period-$t$ action value conditional on the user
choosing Seller A:
\begin{equation}
    Q_{t,x}(S,F)
    \triangleq
    R-c_x
    +p_xV_{t+1}(S+1,F)
    +(1-p_x)V_{t+1}(S,F+1).
    \label{eq:finite-action-value}
\end{equation}
The product choice is made by comparing $Q_{t,1}$ and $Q_{t,2}$. Define the
marginal reputation value of a success relative to a failure by
\begin{equation}
    M_t(S,F)
    \triangleq
    V_{t+1}(S+1,F)-V_{t+1}(S,F+1).
    \label{eq:marginal-reputation}
\end{equation}
Then
\begin{equation}
    Q_{t,2}(S,F)-Q_{t,1}(S,F)
    =
    -(c_2-c_1)+(p_2-p_1)M_t(S,F).
    \label{eq:q-gap}
\end{equation}
Therefore, product 2 is optimal at period $t$ and state $(S,F)$ if and only if
\begin{equation}
    c_2-c_1
    \leq
    (p_2-p_1)M_t(S,F).
    \label{eq:finite-threshold}
\end{equation}
Equivalently,
\begin{equation}
    M_t(S,F)
    \geq
    \frac{c_2-c_1}{p_2-p_1}.
    \label{eq:marginal-threshold}
\end{equation}
The demand probability $D(S,F)$ scales the expected value of acting in the
calendar period, while the product choice itself is governed by the value of
improving the user's future belief. A good signal is valuable when it increases
future demand enough to justify the incremental cost of product 2.

\section{Fluid Approximation}

The exact dynamic program retains the full distribution of the posterior count
process. We next formulate a deterministic fluid approximation that replaces
the random count increments by their conditional expected drifts. This
approximation is useful for identifying the economic forces behind the policy
and for studying long horizons without enumerating the triangular state space.
It does not change the user's information: the user continues to infer Seller
A's quality only from observed successes and failures and never observes the
seller's hidden product choice.

\subsection{Deterministic Count Dynamics}

Let $\tau\in[0,T]$ denote continuous calendar time, and let $s(\tau)$ and
$f(\tau)$ be the fluid success and failure counts. We extend the demand function
in \cref{eq:demand-probability} to nonnegative real counts:
\begin{equation}
    \mathcal{D}(s,f)
    \triangleq
    \int_{p_0}^{1}
    \frac{\theta^s(1-\theta)^f}{B(1+s,1+f)}\,d\theta.
    \label{eq:fluid-demand}
\end{equation}
Let $a(\tau)\in[0,1]$ denote the relaxed intensity of product 2. It can be
interpreted as the fraction of otherwise identical fluid units assigned to
product 2, or as time mixing between the two products. Define
\begin{align}
    q(a) &\triangleq p_1+a(p_2-p_1), \label{eq:fluid-quality}\\
    c(a) &\triangleq c_1+a(c_2-c_1). \label{eq:fluid-cost}
\end{align}
The fluid state then evolves according to
\begin{align}
    \dot{s}(\tau)
    &=
    \mathcal{D}(s(\tau),f(\tau))q(a(\tau)),
    \label{eq:fluid-success-dynamics}\\
    \dot{f}(\tau)
    &=
    \mathcal{D}(s(\tau),f(\tau))[1-q(a(\tau))],
    \label{eq:fluid-failure-dynamics}
\end{align}
with $s(0)=f(0)=0$. The factor $\mathcal{D}(s,f)$ is the fluid rate at which
Seller A is selected and consequently receives a new observation. Conditional
on selection, $q(a)$ and $1-q(a)$ are the expected success and failure rates.

For an admissible measurable control $a(\cdot)$, the fluid profit is
\begin{equation}
    J^{\mathrm{fl}}(a)
    =
    \int_0^T
    \mathcal{D}(s(\tau),f(\tau))
    [R-c(a(\tau))]\,d\tau.
    \label{eq:fluid-objective}
\end{equation}
The fluid control problem is to maximize \cref{eq:fluid-objective} subject to
\cref{eq:fluid-success-dynamics,eq:fluid-failure-dynamics}. This is a
deterministic mean-drift closure: in general,
$\E[D(S_t,F_t)]\neq D(\E[S_t],\E[F_t])$, so the fluid model is an approximation
rather than an exact reformulation of the stochastic problem.

\subsection{Fluid Hamilton--Jacobi--Bellman Equation}

Let $\mathcal{V}(\tau,s,f)$ denote the optimal fluid profit from time $\tau$ to
$T$. Its terminal condition is
\begin{equation}
    \mathcal{V}(T,s,f)=0.
    \label{eq:fluid-terminal}
\end{equation}
The Hamilton--Jacobi--Bellman equation is
\begin{equation}
\begin{aligned}
    -\partial_\tau\mathcal{V}(\tau,s,f)
    =
    \mathcal{D}(s,f)
    \max_{a\in[0,1]}
    \big\{
        &R-c(a)
        +q(a)\partial_s\mathcal{V}(\tau,s,f)\\
        &+[1-q(a)]\partial_f\mathcal{V}(\tau,s,f)
    \big\}.
\end{aligned}
\label{eq:fluid-hjb}
\end{equation}
Write $\Delta p=p_2-p_1>0$ and $\Delta c=c_2-c_1>0$. The terms in the
Hamiltonian that depend on $a$ are
\begin{equation}
    a\left[
        -\Delta c
        +\Delta p\left(
            \partial_s\mathcal{V}(\tau,s,f)
            -\partial_f\mathcal{V}(\tau,s,f)
        \right)
    \right].
    \label{eq:fluid-control-coefficient}
\end{equation}
Consequently, although the fluid control allows mixing, an optimal control is
bang-bang except at points of indifference:
\begin{equation}
    a^*(\tau,s,f)
    =
    \begin{cases}
        1,
        & \Delta p(\partial_s\mathcal{V}-\partial_f\mathcal{V})
          >\Delta c,\\
        0,
        & \Delta p(\partial_s\mathcal{V}-\partial_f\mathcal{V})
          <\Delta c,\\
        \text{any value in }[0,1],
        & \Delta p(\partial_s\mathcal{V}-\partial_f\mathcal{V})
          =\Delta c.
    \end{cases}
    \label{eq:fluid-bang-bang}
\end{equation}
Thus the fluid model preserves the threshold structure of
\cref{eq:finite-threshold}: product 2 is used when the fluid shadow value of
replacing a failure by a success is at least $\Delta c/\Delta p$.

\subsection{Effective Observations and Posterior Mean}

The fluid dynamics become especially transparent after transforming the state.
Let
\begin{equation}
    n(\tau)=s(\tau)+f(\tau),
    \qquad
    m(\tau)=\frac{s(\tau)+1}{n(\tau)+2}.
    \label{eq:fluid-state-transform}
\end{equation}
Here $n$ is the effective number of observations from Seller A and $m$ is the
posterior mean implied by the fluid counts. Define
\begin{equation}
    \widehat{\mathcal{D}}(n,m)
    =
    \Pb\left(
        \operatorname{Beta}((n+2)m,(n+2)(1-m))\geq p_0
    \right).
    \label{eq:fluid-demand-nm}
\end{equation}
Using \cref{eq:fluid-success-dynamics,eq:fluid-failure-dynamics}, the transformed
dynamics are
\begin{align}
    \dot{n}(\tau)
    &=
    \widehat{\mathcal{D}}(n(\tau),m(\tau)),
    \label{eq:fluid-count-dynamics}\\
    \dot{m}(\tau)
    &=
    \frac{\widehat{\mathcal{D}}(n(\tau),m(\tau))}{n(\tau)+2}
    [q(a(\tau))-m(\tau)].
    \label{eq:fluid-mean-dynamics}
\end{align}
Equation \eqref{eq:fluid-mean-dynamics} isolates two forces. The posterior mean
moves toward the quality $q(a)$ delivered by the seller, while its speed of
adjustment declines as $1/(n+2)$. Demand controls the speed of calendar time at
which learning occurs but not the direction of belief updating.

If $\mathcal{W}(\tau,n,m)$ denotes the value function in the transformed state,
then
\begin{equation}
    \partial_s\mathcal{V}-\partial_f\mathcal{V}
    =
    \frac{1}{n+2}\partial_m\mathcal{W}.
    \label{eq:fluid-shadow-transform}
\end{equation}
The fluid product-2 rule can therefore be written as
\begin{equation}
    \frac{\Delta p}{n+2}\partial_m\mathcal{W}(\tau,n,m)
    \geq \Delta c.
    \label{eq:fluid-mean-threshold}
\end{equation}
This representation makes belief rigidity explicit. As the number of
observations grows, one additional success has a smaller effect on the posterior
mean, so the shadow value $\partial_m\mathcal{W}$ must be correspondingly larger
to justify product 2.

\subsection{Long-Run Quality-Mix Benchmark}

Suppose the relaxed control is held constant at $a$ and demand remains positive.
Dividing \cref{eq:fluid-mean-dynamics} by \cref{eq:fluid-count-dynamics} yields
\begin{equation}
    \frac{dm}{dn}=\frac{q(a)-m}{n+2}.
    \label{eq:fluid-mean-by-count}
\end{equation}
Starting from the uniform prior, $n(0)=0$ and $m(0)=1/2$, the solution is
\begin{equation}
    m(n)
    =
    q(a)+\frac{1-2q(a)}{n+2}.
    \label{eq:fluid-mean-solution}
\end{equation}
Hence the posterior mean converges to the average quality $q(a)$. When
$p_1<p_0\leq p_2$, the boundary product-2 intensity that keeps this limiting
mean at the outside-option quality is
\begin{equation}
    a_{\mathrm{bdry}}(p_0)
    =
    \frac{p_0-p_1}{p_2-p_1}.
    \label{eq:fluid-boundary-mix}
\end{equation}
If $p_0\leq p_1$, product 1 alone has quality at least as high as the outside
option, so the analogous boundary intensity is zero. If $p_0>p_2$, no feasible
quality mix can prevent the posterior mean from eventually falling below
$p_0$.

The boundary in \cref{eq:fluid-boundary-mix} is a sustainability benchmark, not
by itself an optimal stationary policy. Posterior concentration implies
\begin{equation}
    \widehat{\mathcal{D}}(n,m(n))
    \longrightarrow
    \begin{cases}
        1, & q(a)>p_0,\\
        1/2, & q(a)=p_0,\\
        0, & q(a)<p_0.
    \end{cases}
    \label{eq:fluid-demand-limit}
\end{equation}
Thus a control strictly above $a_{\mathrm{bdry}}(p_0)$ is needed for demand to
converge to one when $p_1<p_0<p_2$, whereas the exact boundary produces a
singular regime in which the concentrated posterior remains centered at the
user's cutoff. For the computational values $p_1=0.35$ and $p_2=0.80$, the
boundary intensities are $1/3$ at $p_0=0.50$ and $7/9$ at $p_0=0.70$.
Product 1 is already sufficient at $p_0\in\{0.10,0.30\}$, while demand cannot
be sustained asymptotically at $p_0=0.90$ even under product 2.

\subsection{Role and Limitations of the Fluid Model}

The fluid approximation provides three useful benchmarks. First, it retains the
same quality-investment threshold as the exact dynamic program. Second, it
separates the rate at which observations arrive from the rate at which each
observation changes the posterior mean. Third, it identifies the minimum
long-run quality mix needed to keep Seller A competitive with the outside
option.

The approximation deliberately suppresses path dispersion. This matters most
near $m=p_0$, where Thompson-sampling demand changes rapidly and
\cref{eq:fluid-demand-limit} is discontinuous in limiting quality. A diffusion
approximation would be needed to retain the order-$\sqrt{n}$ fluctuations around
that boundary. Accordingly, the fluid solution should be evaluated as a
long-horizon benchmark and compared against the exact backward-induction policy
and stochastic simulation rather than treated as an exact substitute.

\section{Exact Backward Induction}

At the beginning of period $t$, the reachable states satisfy
\begin{equation}
    \State_t=\{(S,F)\in\Natu_0^2:S+F\leq t-1\}.
    \label{eq:reachable-state-space}
\end{equation}
Thus the finite-horizon problem can be solved exactly by backward induction over
the triangular state space. Starting from \cref{eq:terminal-condition}, the
algorithm computes $V_t$ from \cref{eq:finite-bellman} for
$t=T,T-1,\ldots,1$. The implementation stores only the current and next value
arrays during the backward pass. It also records early-period policy slices and
the marginal reputation object $M_t(S,F)$ for diagnostic analysis.

The policy diagnostics focus on four state-dependent objects:
\begin{enumerate}
    \item the marginal reputation value $M_t(S_t,F_t)$;
    \item the demand probability $D(S_t,F_t)$;
    \item the cumulative number of Seller A observations $S_t+F_t$;
    \item the posterior mean $(S_t+1)/(S_t+F_t+2)$.
\end{enumerate}
These objects separate the user's propensity to choose Seller A, the rigidity of
the user's belief, the location of the posterior belief, and the value to Seller
A of generating a success rather than a failure.

\section{Interpretation}

The central object in this finite-horizon problem is $M_t(S,F)$. It measures how
much additional continuation value Seller A obtains from a success rather than a
failure at the next belief update. Product 2 is used exactly when this marginal
reputation value is large enough to pay for its incremental cost. The demand
probability $D(S,F)$ determines how often the seller receives the opportunity to
earn revenue and update the state, while $S+F$ and the posterior mean determine
how responsive the user's belief is to future outcomes. This decomposition makes
the finite-horizon policy transparent: the costly product is not chosen because
it raises current demand, but because it can improve the future reputation state
when the reputational return is high enough. The fluid model preserves this
logic through the shadow-value condition
\cref{eq:fluid-mean-threshold} and complements it with a tractable benchmark for
the long-run quality mix.

\end{document}
